\section{The Leximax Universal-Regret core}
\label{sec:core}

Figure~\ref{fig:pipeline} gives the end-to-end view of the rule developed in this
section: from a candidate set, through normalisation and the declared probe family,
to the leximax selection and its certificate.

% LUR end-to-end pipeline schematic
\begin{figure}[t]
\centering
\begin{tikzpicture}[node distance=6mm and 9mm]
\node[stage] (A) {Candidate set\\$A=\{x_1,\dots,x_N\}$\\\scriptsize(non-dominated)};
\node[stage, right=of A] (N) {Normalise\\$r_i=\dfrac{f_i-\ideal_i}{\nadir_i-\ideal_i}$\\\scriptsize to $[0,1]^m$};
\node[accent, right=of N] (Q) {Declared probes\\$\Q=\{q_1,\dots,q_K\}$\\\scriptsize singletons + clusters};
\node[stage, right=of Q] (D) {Disappointment\\matrix $\reg\in[0,1]^{N\times K}$\\\scriptsize $\reg_q(x)$, Eq.~(2)};
\node[stage, below=10mm of D] (S) {Leximax sort\\$\reg^{\downarrow}(x)$\\\scriptsize worst-first};
\node[win, left=14mm of S] (W) {Recommendation\\$x^{\mathrm{LUR}}$ + certificate\\\scriptsize binding probes named};

\draw[flow] (A) -- (N);
\draw[flow] (N) -- (Q);
\draw[flow] (Q) -- (D);
\draw[flow] (D) -- (S);
\draw[flow] (S) -- (W);

\node[font=\scriptsize\itshape, lurslate, below=1mm of A] {inputs};
\node[font=\scriptsize\itshape, lurblue, below=1mm of Q] {declared, auditable};
\node[font=\scriptsize\itshape, lurcrimson, above=1mm of W] {output};
\end{tikzpicture}
\caption{The LUR pipeline. Modelling choices are confined to the \emph{declared}
probe family $\Q$ and the normalisation bounds; the selection itself is an exact,
deterministic leximax over the disappointment matrix, and the output is a single
recommendation accompanied by an auditable regret certificate.}
\label{fig:pipeline}
\end{figure}


\subsection{Problem setting and normalisation}
Let $\X\subseteq\R^n$ be the feasible set and $F:\X\to\R^m$, $F(x)=(f_1(x),\dots,
f_m(x))$, the vector objective.

\begin{center}
\fbox{
\begin{minipage}{0.9\columnwidth}
\textbf{Convention:} All criteria $f_i$ are assumed to be \emph{minimised}. Maximisation objectives must have their sign reversed before normalisation. Consequently, lower normalised values $r_i$ are better, and probes are non-decreasing functions of these costs.
\end{minipage}
}
\end{center}

For
the finite case let $A=\{x_1,\dots,x_N\}\subseteq\X$ be a candidate set, e.g.\ an
\emph{a posteriori} approximation of the efficient set. Before evaluation, Pareto-dominated alternatives are removed from $A$ unless the analyst explicitly requires selection over the raw un-filtered set. For the continuous case (\S\ref{sec:direct}) $A$ is not required. Criteria are normalised to $[0,1]$ by the
positive affine map
\begin{equation}
r_i(x)=\frac{f_i(x)-\ideal_i}{\nadir_i-\ideal_i},\qquad
\ideal_i=\min_{x\in A}f_i(x),\;\; \nadir_i=\max_{x\in A}f_i(x),
\label{eq:norm}
\end{equation}
which attains both endpoints on $A$. A criterion is flagged as numerically degenerate if the range $\nadir_i-\ideal_i$ approaches zero. Degenerate criteria are reported and excluded from non-singleton probe construction unless the analyst supplies external bounds.

For the finite case all extrema ($\ideal$,
$\nadir$, $q^\star$, $q^{\mathrm w}$ below) are taken over the active candidate set $A$;
for the continuous case (\S\ref{sec:direct}) they are taken over $\X$. We use the
upper attainable value over $A$ in place of the true nadir over the efficient
set; the two coincide when $A$ contains the criterion-wise worst efficient points,
and otherwise the former is a conservative upper normalisation bound. The rule is invariant to
positive affine rescaling of the criteria when the bounds are transformed
consistently; it is not invariant to arbitrary monotone transformations (e.g.\
$f_i\mapsto\log f_i$), which change the distances that define regret---so such
transformations are a modelling choice the analyst must make deliberately.

\subsection{Admissible probes}
A \emph{probe} is a scalar function $q:[0,1]^m\to\R$ encoding a rational decision
question. The family $\Q$ is \emph{declared}: it is part of the specification.

\begin{definition}[Admissible probe]
\label{def:probe}
A probe $q$ is \emph{admissible} if (i) it is non-decreasing in each $r_i$
(monotonicity) and (ii) it can be rescaled so that $q(\mathbf 0)=0$, $q(\mathbf
1)=1$ (normalisation).
\end{definition}

Here $\mathbf 0$ and $\mathbf 1$ denote the all-zeros and all-ones vectors. The initial endpoint normalisation ($q(\mathbf 0)=0, q(\mathbf 1)=1$) serves to make values structurally scale-comparable across the heterogeneous family, before the empirical disappointment anchors ($q^\star, q^{\mathrm w}$) re-normalise the observed active range.
Monotonicity gives \emph{weak} Pareto compatibility; strict compatibility needs a
separation property (Definition~\ref{def:sep}). Canonical admissible probes are
the \emph{singletons} $q_i(r)=r_i$; the \emph{grand mean} $q_\mu(r)=\frac1m
\sum_i r_i$ (the utilitarian question); the \emph{grand maximum}
$q_{\max}(r)=\max_i r_i$ (the egalitarian/worst-criterion question); and, for any
coalition $S\subseteq\{1,\dots,m\}$, the coalition mean and maximum. All have
non-negative weights or are maxima, hence are monotone.

\begin{definition}[Separating family]
\label{def:sep}
A probe family $\Q$ \emph{separates} criteria if for every $i$ there is a probe
$q\in\Q$ strictly increasing in $r_i$. Including all $m$ singleton probes
$q_i(r)=r_i$ is sufficient; LUR always does so.
\end{definition}

\subsection{Disappointment and regret}
For each probe $q\in\Q$ let $q^\star=\min_{x\in A}q(r(x))$ and $q^{\mathrm w}=\max_{x\in A}q(r(x))$
(or a user-supplied worst-acceptable bound). Define the active range $R_q = q^{\mathrm w} - q^\star$. If $R_q < \varepsilon_q$ for a small numerical threshold, the probe is degenerate and is excluded from the primary decision vector. Otherwise, the \emph{normalised disappointment} of $x$ under $q$ is
\begin{equation}
\reg_q(x)=\frac{q(r(x))-q^\star}{R_q}\in[0,1],
\label{eq:disap}
\end{equation}
so that $\reg_q(x)=0$ iff $x$ is $q$-optimal and $1$ iff $x$ attains the worst bound.

We use the term \emph{disappointment} for the normalised shortfall of an alternative from the best attainable value of a single given probe, and \emph{regret certificate} for the vector collecting these probe-wise disappointments. This is regret-like but differs from classical pairwise Savage regret unless the probe family is interpreted specifically as the uncertainty set of possible true preference functions.

\subsection{The leximax order}
Collect $\reg(x)=(\reg_{q_1}(x),\dots,\reg_{q_K}(x))$, $K=|\Q|$, and sort it in
\emph{descending} order to obtain $\reg^{\downarrow}(x)=(\reg_{[1]}(x)\ge\dots\ge
\reg_{[K]}(x))$. The \textbf{Leximax Universal-Regret order} is
\begin{equation}
x\preceq_{\mathrm{LUR}}y \iff \reg^{\downarrow}(x)\lex \reg^{\downarrow}(y),
\label{eq:lur}
\end{equation}
i.e.\ $x$ is preferred if its worst disappointment is smaller; on ties, its
second worst; and so on. The \emph{LUR solution} is $x^{\mathrm{LUR}}\in
\argmin^{\mathrm{lex}}_{x}\reg^{\downarrow}(x)$, and its \emph{certificate} is the
labelled, sorted vector $\reg^{\downarrow}(x^{\mathrm{LUR}})$ together with the
probes attaining the largest entries. Figure~\ref{fig:leximax} illustrates the
comparison on a two-alternative example.

% Leximax comparison of two alternatives' sorted disappointment vectors
\begin{figure}[t]
\centering
\begin{tikzpicture}[x=1cm, y=1cm]
\def\barw{0.42}
\def\sc{0.8}        % vertical scale for bar heights
\def\rowA{3.0}      % baseline y for row A
\def\rowB{1.6}      % baseline y for row B
\def\Lx{0}          % left block origin x
\def\Rx{4.8}        % right block origin x

% --- column titles ---
\node[font=\scriptsize\bfseries] at (\Lx+1.1,4.45) {unsorted regrets};
\node[font=\scriptsize\bfseries] at (\Rx+1.1,4.45) {sorted $\reg^{\downarrow}$ (worst-first)};

% --- row labels ---
\node[font=\small, anchor=east] at (\Lx-0.15,\rowA+0.3) {$A$};
\node[font=\small, anchor=east] at (\Lx-0.15,\rowB+0.3) {$B$};
\node[font=\small, anchor=east] at (\Rx-0.15,\rowA+0.3) {$A$};
\node[font=\small, anchor=east] at (\Rx-0.15,\rowB+0.3) {$B$};

% --- unsorted bars ---
\foreach \v/\i in {0.2/0,0.9/1,0.4/2,0.7/3,0.3/4}{
  \fill[lexurslate!50] (\Lx+\i*0.5,\rowA) rectangle ++(\barw,{\v*\sc});}
\foreach \v/\i in {0.6/0,0.8/1,0.5/2,0.8/3,0.2/4}{
  \fill[lexurslate!50] (\Lx+\i*0.5,\rowB) rectangle ++(\barw,{\v*\sc});}

% --- sorted bars (A wins/loses comparison) ---
\foreach \v/\i in {0.9/0,0.7/1,0.4/2,0.3/3,0.2/4}{
  \fill[lexurslate!50] (\Rx+\i*0.5,\rowA) rectangle ++(\barw,{\v*\sc});}
\foreach \v/\i in {0.8/0,0.8/1,0.6/2,0.5/3,0.2/4}{
  \fill[lexurslate!50] (\Rx+\i*0.5,\rowB) rectangle ++(\barw,{\v*\sc});}
% highlight the deciding first coordinate
\draw[lexurcrimson, line width=1pt] (\Rx,\rowA) rectangle ++(\barw,{0.9*\sc});
\node[lexurcrimson, font=\scriptsize] at (\Rx+\barw/2,\rowA+0.9*\sc+0.18) {$0.9$};
\draw[lexurgreen!60!black, line width=1pt] (\Rx,\rowB) rectangle ++(\barw,{0.8*\sc});
\node[lexurgreen!60!black, font=\scriptsize] at (\Rx+\barw/2,\rowB+0.8*\sc+0.18) {$0.8$};

% --- comparison note (to the right, clear of bars) ---
\node[anchor=west, font=\scriptsize, align=left, lexurslate] at (\Rx+2.9,2.3)
  {compare worst-first:\\[1pt]
   $0.8<0.9 \;\Rightarrow\; B\prec_{\mathrm{LexUR}}A$};
\draw[flow, draw=lexurcrimson] (\Rx+2.85,2.3) -- (\Rx+0.5,2.5);

% --- verdict box ---
\node[win, font=\scriptsize, align=center] at (\Rx+0.0,0.45)
  {$B$ is recommended: its \emph{worst} disappointment is smallest.
   On a tie at the worst entry, the next sorted coordinate decides, and so on.};
\end{tikzpicture}
\caption{The leximax order. Each alternative's probe disappointments are sorted
worst-first; alternatives are then compared lexicographically on the sorted
vectors. Here $B$ wins because its largest disappointment ($0.8$) is below $A$'s
($0.9$); had they tied, the second-largest entries would decide, and so on.}
\label{fig:leximax}
\end{figure}


Exact LUR defines the mathematical complete preorder. In computation and practical reporting, a tolerance $\tau>0$ is used only to define a numerical indifference class: we report alternatives whose sorted disappointment vectors are not separated from the exact winner by more than $\tau$ at the first differing coordinate. The $\tau$-tolerance rule is not itself a transitive preorder, but rather a numerical reporting and stability device around the exact or margin-separated winner (Theorem~\ref{thm:stab}).

\subsection{Foundational guarantees}
\begin{theorem}[Existence and completeness of exact LUR]
\label{thm:exist}
For any finite $A$ and finite $\Q$, $x^{\mathrm{LUR}}$ exists and exact
$\preceq_{\mathrm{LUR}}$ is a complete preorder on $A$.
\end{theorem}
\begin{proof}
For finite $A$ and finite $\Q$, each $\reg^{\downarrow}(x)\in[0,1]^K$ is a
well-defined real vector. The lexicographic order on $\R^K$ is reflexive,
transitive and total, hence a complete preorder; since $A$ is finite the
lexicographic minimum exists.
\end{proof}

\begin{theorem}[Weak and strict Pareto compatibility]
\label{thm:pareto}
Let $x$ strictly Pareto-dominate $y$ and let every $q\in\Q$ be monotone
non-decreasing in each $r_i$. Then $\reg_q(x)\le\reg_q(y)$ for all $q$
(\emph{weak} compatibility), so $x\preceq_{\mathrm{LUR}}y$. If, in addition, $\Q$
\emph{separates} criteria (Definition~\ref{def:sep})---in particular if $\Q$
contains the singleton probes---then $\reg_q(x)<\reg_q(y)$ for some $q$ and
$x\prec_{\mathrm{LUR}}y$.
\end{theorem}
\begin{proof}
Let $x\prec_P y$, so $r_i(x)\le r_i(y)$ for all $i$ with strict inequality for
some $j$. For any monotone non-decreasing probe $q$, $q(r(x))\le q(r(y))$; as
$q^\star,q^{\mathrm w}$ are constants independent of the alternative,
$\reg_q(x)\le\reg_q(y)$ for all $q$, giving $\reg^{\downarrow}(x)\le
\reg^{\downarrow}(y)$ componentwise after sorting and hence
$x\preceq_{\mathrm{LUR}}y$ (weak compatibility, monotonicity only). If $\Q$
separates criteria there is a probe $q$ strictly increasing in $r_j$, for which
$q(r(x))<q(r(y))$ and thus $\reg_q(x)<\reg_q(y)$; the multiset of disappointments
of $x$ is then componentwise $\le$ that of $y$ with at least one strict drop, so
the descending-sorted vectors satisfy $\reg^{\downarrow}(x)\lex\reg^{\downarrow}
(y)$ strictly and $x\prec_{\mathrm{LUR}}y$. The counterexample in the main text
shows the separation hypothesis cannot be dropped.
\end{proof}

The separation hypothesis is necessary: with only the probe $q(r)=r_2$, the
dominator $r(x)=(0,1)$ and the dominated $r(y)=(0.1,1)$ have equal disappointment,
so monotonicity alone does not give a strict preference. LUR always includes the
singletons, so it separates.

\begin{corollary}[Dominated exclusion]
\label{cor:dom}
If $x$ is Pareto-dominated by a feasible $y$ and $\Q$ separates criteria, then
$x$ is not LUR-optimal.
\end{corollary}
\begin{proof}
If $x$ is dominated by a feasible $y$ and $\Q$ separates criteria,
Theorem~\ref{thm:pareto} gives $y\prec_{\mathrm{LUR}}x$, so $x$ is not a
lexicographic minimiser. Without separation only $y\preceq_{\mathrm{LUR}}x$ holds
and a dominated alternative may tie its dominator.
\end{proof}

\begin{theorem}[Stability of the implemented winner under a margin condition]
\label{thm:stab}
Let $\hat\reg_q(x)$ estimate $\reg_q(x)$
with $\eta=\max_{x,q}|\hat\reg_q(x)-\reg_q(x)|$. Suppose the exact winner
$x^{\mathrm{LUR}}$ beats every rival $y$ at the first coordinate $k(y)$ at which
their sorted vectors differ by more than $\tau$, with margin
$\reg_{[k(y)]}(y)-\reg_{[k(y)]}(x^{\mathrm{LUR}})\ge\Delta$, and that all earlier
coordinates agree within $\tau-2\eta$. If $2\eta<\min\{\Delta-\tau,\ \tau\}$, the
winner under $\hat\reg$ is unchanged.
\end{theorem}
\begin{proof}
Write $\hat\reg_q(x)=\reg_q(x)+\xi_q(x)$, $|\xi_q(x)|\le\eta$; sorting is
$1$-Lipschitz in $\|\cdot\|_\infty$, so each sorted coordinate moves by at most
$\eta$. Fix a rival $y$ and the deciding coordinate $k=k(y)$ with true margin
$\reg_{[k]}(y)-\reg_{[k]}(x^{\mathrm{LUR}})\ge\Delta$. Under perturbation this
margin is at least $\Delta-2\eta$; since $2\eta<\Delta-\tau$ it stays above
$\tau$, so the $\tau$-tolerance rule still decides coordinate $k$ in favour of
$x^{\mathrm{LUR}}$. For the earlier coordinates $j<k$, which agree within
$\tau-2\eta$ by hypothesis, the perturbed gaps stay within $\tau$ and are treated
as ties, so the decision is not pre-empted before $k$. As $2\eta<\tau$ no spurious
strict difference is created earlier. Hence $x^{\mathrm{LUR}}$ still beats every
rival and the winner is unchanged. (Exact comparison, $\tau=0$, is excluded
precisely because the second condition $2\eta<\tau$ then fails.)
\end{proof}

Theorem~\ref{thm:stab} is deliberately stated for the tolerance rule: exact
lexicographic comparison is \emph{not} perturbation-stable, because a winner that
beats a rival only at a later coordinate after an exact earlier tie can be
reversed by an arbitrarily small perturbation of the earlier coordinate. The
tolerance $\tau$ absorbs exactly those near-ties. We
note that the margin condition $\Delta$ in Theorem~\ref{thm:stab} can be
restrictive on densely packed candidate sets, where the separation between
certificates may be very small. Without that margin the theorem gives no winner
stability guarantee; the appropriate output is therefore a tolerance or
stability set rather than a uniquely stable recommendation. For
continuous $\X$, existence requires compactness of $\X$ and continuity of each
$q\circ r$; uniqueness is not guaranteed and a small indifference class may arise.

